\documentclass[11pt,a4paper]{article}
\usepackage[T1]{fontenc}
\usepackage[utf8]{inputenc}
\usepackage[french]{babel}
\usepackage{lmodern}
\usepackage{amsmath,amssymb}
\usepackage[margin=2.5cm]{geometry}
\usepackage{booktabs}
\usepackage{graphicx}
\usepackage{enumitem}
\usepackage{xcolor}
\usepackage{tabularx}
\usepackage{hyperref}
\hypersetup{colorlinks=true,linkcolor=blue!50!black}

\newcommand{\NW}{\mathrm{NW}}
\newcommand{\code}[1]{\texttt{#1}}
\newcommand{\Workspace}{\code{m4\_credit\_soc\_fable}}
\newcommand{\OtherWorkspace}{\code{m4\_credit\_soc\_sol}}

\title{Prompt de recherche M4 — avalanches auto-organisées\\
\large Version pour Claude Fable~5 (workspace \Workspace{})}
\author{Anatole Joseph Stouls}
\date{13 juillet 2026}

\begin{document}
\maketitle

\begin{abstract}
Fork de M3 (\code{m3\_credit\_soc}) : société de crédit multi-agents en
recherche d'un régime d'avalanches de faillites auto-organisé (SOC), robuste
en loi de puissance, dépassant nettement les tailles obtenues jusqu'ici
(max$=7$), avec un code aussi simple que possible. Ce document est le brief
de recherche remis à Claude Fable~5 pour cette phase~; les instructions
d'exécution propres au modèle sont en fin de document.
\end{abstract}

\section*{Lecture préalable obligatoire}

\begin{itemize}[leftmargin=*]
  \item \code{m3\_credit\_soc/NOTES.md} — journal complet du programme M3.
    Utile comme historique de ce qui a déjà été essayé, \textbf{pas comme
    verdict à respecter} (voir plus bas). Ne relancez pas un test déjà fait
    sans le citer et dire pourquoi le refaire apporte quelque chose de
    nouveau.
  \item \code{m3\_credit\_soc/reports/05\_negative\_results/main.tex} et
    \code{reports/07\_income\_rule/main.tex} — synthèses rédigées des
    résultats négatifs et du mini-protocole X1, même remarque.
  \item \code{src/m4/*.py} (dans \Workspace{}) — moteur actuel, fork de M3.
    \code{config.py} documente chaque champ et pourquoi ses valeurs par
    défaut sont ce qu'elles sont.
  \item \textbf{Articles et notes de conception de référence, pour les lois
    réelles et les pistes de mécanisme} — dossier
    \code{recherche/} (racine du dépôt \code{jupyter/})~:
    \begin{itemize}
      \item \code{recherche/analyse\_articles/} — recodage et analyse de
        modèles d'économie statistique~: Bouchaud~\&~Mézard, \emph{Wealth
        condensation in a simple model of economy} (mécanisme de
        condensation de richesse par croissance multiplicative + échange,
        directement pertinent pour une dynamique d'accumulation-relaxation)~;
        deux modèles de Ian Wright sur les microfondations et l'architecture
        sociale d'une économie de classes. Le texte intégral des articles
        n'est pas versionné (copyright) mais chaque sous-dossier contient
        le recodage, l'analyse statistique et un rapport.
      \item \code{recherche/analyse\_distributions\_taille\_revenu/} —
        boîte à outils déjà écrite pour tester la famille de loi (13
        familles, MLE, AIC/BIC) et la queue (test de Clauset--Shalizi--Newman)
        de distributions de taille/revenu — à réutiliser, pas à réinventer.
      \item \code{recherche/conception\_boltzmann\_pareto\_soc/} — note de
        conception antérieure à M2, qui discute déjà explicitement une
        transition stable~$\to$~SOC pilotée par un paramètre de connectivité
        du marché et des diagnostics SOC (tailles de cascade, corrélation
        concentration/défaut). Directement pertinent pour la conception du
        mécanisme recherché ici.
    \end{itemize}
\end{itemize}

\section*{Contexte}

M3 est un modèle multi-agents de société de crédit : entités avec liquidité
\(L\) et capital productif \(K\), prêts bilatéraux, faillites en cascade,
recherche d'un régime SOC (\emph{self-organized criticality}) sur la taille
des avalanches de faillites. M4 est un fork de M3. Un premier run simple
(\code{m4\_first\_s0}, seed~0, $T{=}2000$, baseline chocs sectoriels
corrélés + objectif de revenu myope) a été exécuté et donne une population
effondrée ($\approx 135$) et des avalanches non triviales mais petites
(max$=7$).

\textbf{Vous travaillez dans \Workspace{}, un dossier de travail
indépendant.} Un autre agent travaille en parallèle dans \OtherWorkspace{},
sur le même sujet, à partir du même point de départ. Les deux dossiers sont
des copies indépendantes~: ne modifiez pas \OtherWorkspace{}. Vous pouvez
consulter le \code{NOTES.md} de l'autre dossier s'il existe et si cela vous
aide, mais ne convergez pas prématurément vers la même solution sous ce
seul prétexte — l'intérêt de deux dossiers séparés est d'obtenir deux
explorations réellement indépendantes.

\textbf{Avertissement fort~: les résultats de M2 et de M3 sont à prendre
avec des pincettes.} Ils viennent d'un programme antérieur, sous des choix
de conception et des critères qui ne sont pas forcément les bons pour cette
phase (voir en particulier, plus bas, la correction apportée sur l'effet du
réseau de crédit, qui avait été mal généralisé dans une version précédente
de ce document). Traitez tout ce qui suit comme un point de départ
discutable, pas comme un acquis.

\section*{Objectif}

Obtenir un régime où les avalanches de faillites (composantes connexes du
graphe de pertes, cf. \code{bankruptcy.py::\_build\_avalanches}) sont
produites par un mécanisme d'\textbf{accumulation-relaxation endogène} — pas
par amplification d'une corrélation imposée de l'extérieur — et suivent une
distribution de taille en loi de puissance robuste, avec des tailles
significativement au-delà de 7 (potentiellement 100+ selon la taille du
système). Le système doit continuer à produire des distributions de
\textbf{taille des entités} ($\NW$, $K$ — pas la taille des avalanches) et de
\textbf{revenu} dont la \textbf{famille} de loi colle aux lois réelles
étudiées dans \code{recherche/} (voir « Contraintes à préserver »).

\textbf{Objectif final de simplification~: un code simple.} La liquidité
d'illiquidité (défaut d'une entité solvable mais illiquide) \emph{n'est pas
un enjeu de ce programme} — ne cherchez pas à préserver ce canal pour
lui-même. Simplifier les variables d'état (par exemple fusionner $L$ et $K$,
ou toute autre réduction) est encouragé si ça allège le code et que les
contraintes ci-dessous restent respectées. En particulier, $d_0$ doit à
terme être \textbf{complètement retiré du code} (pas seulement fixé à~0,
voir plus bas) — c'est un objectif de simplification explicite, pas une
option parmi d'autres.

\paragraph{Refonte radicale autorisée.} Vous n'êtes pas limité·e à ajuster
des paramètres de M4 tel quel. Toute partie du mécanisme peut être repensée
si l'analyse le justifie et si le changement est documenté et comparé à la
baseline actuelle~: la règle de faillite (\code{bankruptcy.py}), la
topologie de marché (\code{market.py}), la structure du choc, le nombre et
la nature des variables d'état par entité, la règle de taux, etc. Rien
n'est sanctuarisé sauf ce qui est listé plus bas comme contrainte.

\paragraph{Point de départ retenu pour l'investigation~: la sectorialité
des chocs.} Commencez par l'axe des chocs sectoriels corrélés
(\code{shock\_rho\_sector}, \code{n\_sectors} — variantes G/G2 de M3), pas
par la piste de dette de subsistance (qui reste une voie possible mais pas
privilégiée, voir plus bas). C'est un choix délibéré~: M3 a lu cet axe comme
une amplification continue sans transition, mais ce jugement reposait
notamment sur le critère var/mean, dont ce document explique plus bas
pourquoi il n'est pas nécessaire — l'axe mérite d'être ré-examiné avec le
bon critère.

\paragraph{Marqueur déjà observé, à prendre au sérieux~: régression en loi
de puissance après exclusion des avalanches de taille 1.} En excluant les
avalanches de taille 1 de la régression log-log (garder les deux
régressions, avec et sans taille 1, affichées), on obtient déjà de très bons
ajustements ($r^2 > 0{,}90$) sur des runs existants. C'est un très bon
marqueur d'un régime auto-critique — bien meilleur que var/mean — à
calculer systématiquement pour chaque mécanisme candidat (voir « Critères de
succès »).

\section*{Ce que M3 a établi (points de départ, à ré-examiner, pas des
vérités intangibles)}

\begin{enumerate}[leftmargin=*]

\item \textbf{Corrélation imposée vs auto-organisation — axe de départ de
  l'investigation.} G1 (choc macro) et G2/G2b/G2c (choc sectoriel, $\rho_s$
  de 0 à 0,8) montraient dans M3 une réponse jugée \emph{continue}~: var/mean
  des avalanches de 0,07 (iid) à 0,84 ($\rho_s{=}0{,}8$), max jusqu'à 72,
  sans transition nette au sens de ce critère. Mais var/mean n'est pas le
  bon critère (voir « Critères de succès ») — c'est précisément pourquoi cet
  axe est repris en premier ici, avec le marqueur de régression hors
  taille~1 comme signal principal.

\item \textbf{Décision de conception retenue pour M4~: $d_0$ est supprimé.}
  Le run H de M3 ($d_0{=}0$) a montré que le système se borne seul par
  l'insolvabilité contractuelle (population stationnaire, renouvellement
  observé) sans plancher exogène~: on sait donc qu'on peut s'en passer.
  Commencez par \code{d0=0} en configuration si besoin d'une étape
  intermédiaire, mais l'objectif final est de retirer \code{d0} et son
  usage \textbf{du code lui-même} (\code{config.py}, \code{bankruptcy.py::
  net\_worth}), pas seulement de le neutraliser par sa valeur. La piste
  consistant à router une dette de subsistance vers de vraies créancières du
  marché (plutôt que vers un puits abstrait) reste une voie possible à
  explorer, mais ce n'est \textbf{pas la seule ni la prioritaire} — l'axe
  retenu en premier est la sectorialité des chocs (ci-dessus).

\item \textbf{Le réseau de crédit a un effet réel et attendu sur les
  cascades.} Activer le crédit (vs l'ablation B, sans marché) multiplie les
  cascades~: c'est le canal même qui les rend possibles en connectant les
  entités entre elles, et c'est parfaitement normal — ne traitez aucun
  résultat de ce document comme disant le contraire. Ce que la dose-réponse
  F/F''/baseline de M3 a réellement fermé est plus étroit~: \emph{à volume
  de crédit égal}, changer la règle d'appariement (assortative~/ random~/
  random\_lender) ne changeait rien de mesuré au-delà de ce volume — mais ce
  résultat lui-même est à ré-examiner avec les bons critères plutôt qu'à
  prendre pour acquis.

\item \textbf{Le crédit est causal démographiquement, pas
  distributionnellement}, en baseline richesse (A vs B) selon M3~; ce n'est
  qu'avec la règle de revenu (X1) que le crédit devenait distributionnellement
  actif (Gini $0{,}62$--$0{,}69$ vs $0{,}56$). À vérifier plutôt qu'à
  supposer si votre mécanisme candidat change la donne.

\end{enumerate}

\section*{Contraintes à préserver}

\begin{itemize}[leftmargin=*]
  \item \textbf{Familles de distribution.} Corps de $\NW$, $K$, revenu
    doivent rester dans des familles reconnaissables (pas nécessairement
    identiques à Fisk/dPlN/lognormal observées en M3 — à re-tester, pas à
    supposer). Pour les lois réelles de référence, voir
    \code{recherche/analyse\_articles/} et
    \code{recherche/analyse\_distributions\_taille\_revenu/} (lecture
    préalable ci-dessus) plutôt que de partir de zéro.
  \item \textbf{Renouvellement démographique minimal, pas un taux de
    mortalité cible.} La seule exigence est qu'il existe un renouvellement
    (turnover) de la population/du haut de la distribution, même très
    lent — pas de taux de mortalité particulier à atteindre, et ne
    réutilisez pas un critère de mortalité instantanée du top (voir pièges
    d'outillage) : mesurez le renouvellement entre fenêtres temporelles.
  \item \textbf{Invariants comptables du moteur} (cascade de faillite avec
    point fixe garanti, snapshot en lecture pure pour les figures, etc. —
    voir docstrings \code{bankruptcy.py}, \code{contracts.py}, tests
    \code{tests/}). Toute refonte du moteur doit préserver un jeu
    d'invariants équivalent et le tester.
  \item \textbf{Reproductibilité}~: \code{random.Random(seed)} isolé,
    jamais le RNG global~; artefacts sur disque (\code{config.json},
    \code{summary.json}, \code{series.csv}, \code{avalanches.csv},
    \code{snap\_t*.npz}) comme seule source pour les figures — jamais
    l'état mémoire de la simulation.
\end{itemize}

\section*{Pièges d'outillage à ne pas reproduire}

\begin{enumerate}[leftmargin=*]
  \item Test LR Pareto/lognormale \textbf{sans renormalisation} de la
    lognormale tronquée $\rightarrow$ biais pro-Pareto systématique.
  \item AIC du corps de distribution avec des familles \textbf{non
    tronquées} sur un corps qui est en réalité tronqué $\rightarrow$ biais
    anti-exponentiel.
  \item Transfert prorata des créances de faillite sans fusion par paire
    (\code{merge\_pairs}) $\rightarrow$ explosion combinatoire du carnet de
    prêts (observé jusqu'à 12,4~M de contrats sur la grille M3).
  \item Moyenne arithmétique des taux d'intérêt au lieu de géométrique
    $\rightarrow$ le marché de crédit s'éteint de lui-même.
  \item Critère de « mortalité instantanée du top » — mauvais critère,
    structurellement nul. Ce qu'on veut n'est qu'un \textbf{renouvellement},
    même très lent~: mesurez-le entre fenêtres temporelles, pas
    instantanément.
  \item \textbf{Période transitoire.} Ne calculez pas de statistiques sur
    des données qui incluent encore le régime transitoire de démarrage —
    identifiez et excluez le burn-in avant d'ajuster une distribution ou de
    compter des avalanches.
  \item \textbf{Robustesse de la fenêtre d'analyse.} Faites varier le temps
    de simulation $T$ pour vérifier que le choix de fenêtre d'analyse (où
    commence le régime stationnaire, sur quelle plage on ajuste) est
    robuste, pas un artefact d'un $T$ particulier.
  \item \textbf{Variance inter-seeds.} Pour chaque point de paramètres,
    utilisez plusieurs seeds afin d'estimer la variance des résultats — un
    seed unique ne suffit pas pour juger qu'un effet est réel.
\end{enumerate}

\section*{Critères de succès}

\begin{enumerate}[leftmargin=*]
  \item \textbf{Marqueur principal~: régression en loi de puissance après
    exclusion des avalanches de taille 1, $r^2 > 0{,}90$.} Gardez les deux
    régressions affichées (toutes tailles, et hors taille~1) — c'est le
    signal le plus net déjà observé d'un régime auto-critique. À calculer
    systématiquement pour chaque mécanisme candidat.
  \item \textbf{Scaling en taille finie}~: le cutoff de la distribution de
    taille d'avalanches croît avec la taille du système (population
    stationnaire ou nombre d'entités actives), signature classique d'un
    régime critique.
  \item \textbf{Ajustement en loi de puissance qui survit au test LR
    corrigé} (avec la renormalisation correcte des alternatives tronquées —
    piège n°1 ci-dessus), pas seulement une droite à l'œil sur un graphe
    log-log.
  \item \textbf{La condition sur-Poisson (var/mean $>1$) N'EST PAS
    nécessaire.} Ne l'utilisez pas comme critère de rejet d'un mécanisme —
    M3 en a fait un usage excessif alors que ce n'est pas ce qui compte
    (voir le marqueur principal ci-dessus).
  \item Les distributions de taille des entités et de revenu gardent une
    famille reconnaissable, en cohérence avec les lois réelles de
    référence (\code{recherche/}), avec le changement documenté si une
    famille change.
  \item Le mécanisme proposé a une \textbf{histoire causale endogène}
    articulable en une phrase (pas juste « on a poussé un paramètre de
    corrélation plus haut ») — feedback d'accumulation puis de relaxation
    porté par la dynamique du modèle lui-même.
  \item \textbf{Piège à nommer explicitement~: un effondrement quasi total
    n'est pas de la SOC.} La façon la moins chère d'obtenir de « grosses
    avalanches » dans un modèle de cascade de levier est un régime bimodal
    (une masse de défauts de taille~1 + un effondrement occasionnel de tout
    le système) — ce n'est ni une loi de puissance ni un régime critique,
    c'est une extinction. Un run qui rapporte « avalanche de 100+
    obtenue~! » doit être vérifié contre ce scénario (et contre les
    critères 1--3) avant d'être compté comme un succès.
\end{enumerate}

\section*{Livrables attendus}

\begin{itemize}[leftmargin=*]
  \item Modifications du moteur dans \code{src/m4/} (ou nouveau module si
    la refonte le justifie), avec tests couvrant les invariants touchés.
  \item Au moins un run comparatif baseline-vs-mécanisme-candidat avec
    figures (réutiliser les recettes déjà établies~:
    \code{simulation\_lab/plot\_utils.py} et \code{webapp/mpl\_figures.py} —
    log-binning adaptatif, erreur bootstrap, double régression avec/sans
    taille~1, écrites pour \code{cascades\_rank\_size.png} dans
    \code{m3\_credit\_soc} — à copier/adapter, pas à réinventer).
  \item Une entrée de journal dans \code{\Workspace{}/NOTES.md} (sur le
    modèle de celui de M3~: un résumé d'une ligne en tête de chaque entrée,
    confirmé ou infirmé, et pourquoi) pour chaque hypothèse testée, y
    compris les échecs — les négatifs sont aussi précieux que les positifs
    dans ce programme.
\end{itemize}

\section*{Quand s'arrêter et demander}

\textbf{Ne vous arrêtez pas pour demander confirmation.} Continuez le
travail de bout en bout, y compris pour les décisions notables (changer une
convention documentée ailleurs comme figée, lancer un run coûteux, conclure
qu'aucun mécanisme testé ne produit de SOC robuste). Dans tous ces cas,
\textbf{notifiez-le dans le journal / le rapport} — ce qui compte est que la
décision et sa justification soient tracées, pas qu'elle soit validée avant
d'être prise.



\section*{Instructions d'exécution — Claude Fable 5}

Vous êtes l'agent de développement chargé de cette phase de recherche.
Quelques attentes propres à la façon dont vous travaillez sur des tâches
longues et ouvertes~:

\begin{itemize}[leftmargin=*]
  \item \textbf{Mémoire de travail.} Tenez \code{m4\_credit\_soc/NOTES.md}
    comme journal au fil de l'eau, une entrée par hypothèse testée (cf.
    « Livrables attendus »)~: résumé d'une ligne en tête, confirmé/infirmé,
    et pourquoi. Ne dupliquez pas ce que l'historique git ou le code
    documente déjà.
  \item \textbf{Vérification avant de rapporter.} Avant d'annoncer un
    résultat comme acquis, auditez chaque affirmation contre une sortie
    d'outil réelle de cette session (run, test, figure). Si quelque chose
    n'est pas vérifié, dites-le explicitement plutôt que de l'affirmer.
  \item \textbf{Subagents pour la vérification longue.} Pour les runs
    longs ou les ajustements statistiques (tests LR, scaling en taille
    finie), préférez un subagent à contexte frais pour re-vérifier plutôt
    que l'auto-critique seule.
  \item \textbf{Ambition.} Ce sujet est délibérément ouvert (refonte
    radicale autorisée) — c'est le type de tâche où vous performez le
    mieux. N'hésitez pas à proposer et comparer plusieurs mécanismes
    candidats plutôt que de vous limiter au premier qui fonctionne un peu.
  \item \textbf{Autonomie.} Ne vous arrêtez pas pour demander confirmation,
    même sur les décisions notables — notifiez-les dans le journal (voir
    « Quand s'arrêter et demander » ci-dessus) et continuez. Ne vous arrêtez
    que sur une action destructive/irréversible/hors-périmètre.
  \item \textbf{Dossier de travail indépendant.} Travaillez exclusivement
    dans \Workspace{}. Ne modifiez pas \OtherWorkspace{} (l'autre agent y
    travaille en parallèle sur le même sujet, indépendamment).
\end{itemize}

\end{document}
